% GENERATED FILE. Edit canonical lessons, then run npm run generate:books.
% canonical-source-hash: fnv1a64:024e329eb4cde247
% canonical-lessons: BN-C171-kuri, BN-C171-golap, BN-C171-bot, BN-C171-dhan, BN-C171-joyar

\chapter{Bud, Rose, Banyan tree, Paddy, High tide}
\label{ch:bn-bud-rose-banyan-tree-paddy-high-tide}
\hlchaptermodality{171}

\begin{chapteropening}
\textbf{By the end of this chapter:} \emph{I can say a bud, a rose, a banyan tree, paddy and high tide.}

Complete the last lesson of chapter 171: I can say a bud, a rose, a banyan tree, paddy and high tide.
\end{chapteropening}

\section[\texorpdfstring{\bn{কুঁড়ি} (kũṛi)}{কুঁড়ি (kũṛi)}]{\bn{কুঁড়ি} (kũṛi) — a bud}

\label{lesson:BN-C171-kuri}



\begin{quote}
Before the new one: say the Bengali for a procession, then the Bengali for a tap.
\end{quote}

\subsection*{You'll want to know: \bn{কুঁড়ি}}
\textbf{\bn{কুঁড়ি}} — \emph{kũṛi} — \textquotedblleft{}a bud\textquotedblright{}.

\begin{cousinweb}
The origin of \textbf{\bn{কুঁড়ি}} (\emph{kũṛi}), \textquotedblleft{}a bud\textquotedblright{}, is not settled. CDIAL lists a Bengali \emph{kũṛi} only as \textquotedblleft{}earthen pot\textquotedblright{} (CDIAL 3264), a different word, and a search by the gloss \textquotedblleft{}bud\textquotedblright{} finds only \emph{kali} (CDIAL 2934) and \emph{bol} (CDIAL 10146).
\end{cousinweb}

\begin{grammarlens}[title={What you've built}]
One more word for this chapter.
\end{grammarlens}

\subsection*{Your turn}
\begin{itemize}
\raggedright
  \item \emph{Say it:} \emph{kũṛi}
  \item \emph{Say it:} \emph{kũṛi}, once more
  \item \emph{Say it:} say \emph{kũṛi}
  \item \emph{Recall:} say the Bengali for a procession, then the Bengali for a tap, then say \emph{kũṛi} again
\end{itemize}

\begin{tcolorbox}[breakable,colback=teal!4,colframe=teal!35!black,title={Before you move on}]
What does \bn{কুঁড়ি} mean? (\textquotedblleft{}A bud\textquotedblright{}.) Say it once more.
\end{tcolorbox}
\section[\texorpdfstring{\bn{গোলাপ} (golāp)}{গোলাপ (golāp)}]{\bn{গোলাপ} (golāp) — a rose}

\label{lesson:BN-C171-golap}



\begin{quote}
Before the new one: say the Bengali for a tap, then the Bengali for a bud.
\end{quote}

\subsection*{You'll want to know: \bn{গোলাপ}}
\textbf{\bn{গোলাপ}} — \emph{golāp} — \textquotedblleft{}a rose\textquotedblright{}.

\begin{cousinweb}
\textbf{\bn{গোলাপ}} (\emph{golāp}) is a loan, usually traced to Persian \emph{gulāb}, \textquotedblleft{}rose-water\textquotedblright{}. Persian \emph{gulāb} is \emph{gul}, \textquotedblleft{}rose\textquotedblright{}, + \emph{āb}, \textquotedblleft{}water\textquotedblright{}.
\end{cousinweb}

\begin{grammarlens}[title={What you've built}]
One more word for this chapter.
\end{grammarlens}

\subsection*{Your turn}
\begin{itemize}
\raggedright
  \item \emph{Say it:} \emph{golāp}
  \item \emph{Say it:} \emph{golāp}, once more
  \item \emph{Say it:} say \emph{golāp}
  \item \emph{Recall:} say the Bengali for a tap, then the Bengali for a bud, then say \emph{golāp} again
\end{itemize}

\begin{tcolorbox}[breakable,colback=teal!4,colframe=teal!35!black,title={Before you move on}]
What does \bn{গোলাপ} mean? (\textquotedblleft{}A rose\textquotedblright{}.) Say it once more.
\end{tcolorbox}
\section[\texorpdfstring{\bn{বট} (bôṭ)}{বট (bôṭ)}]{\bn{বট} (bôṭ) — a banyan tree}

\label{lesson:BN-C171-bot}



\begin{quote}
Before the new one: say the Bengali for a bud, then the Bengali for a rose.
\end{quote}

\subsection*{You'll want to know: \bn{বট}}
\textbf{\bn{বট}} — \emph{bôṭ} — \textquotedblleft{}a banyan tree\textquotedblright{}.

\begin{cousinweb}
\textbf{\bn{বট}} (\emph{bôṭ}) is Sanskrit \emph{vaṭa}, \textquotedblleft{}banyan\textquotedblright{}, taken over whole as a learned loan, a tatsama (CDIAL 11211). The word inherited from it is \emph{baṛ}, which CDIAL lists.
\end{cousinweb}

\begin{grammarlens}[title={What you've built}]
One more word for this chapter.
\end{grammarlens}

\subsection*{Your turn}
\begin{itemize}
\raggedright
  \item \emph{Say it:} \emph{bôṭ}
  \item \emph{Say it:} \emph{bôṭ}, once more
  \item \emph{Say it:} say \emph{bôṭ}
  \item \emph{Recall:} say the Bengali for a bud, then the Bengali for a rose, then say \emph{bôṭ} again
\end{itemize}

\begin{tcolorbox}[breakable,colback=teal!4,colframe=teal!35!black,title={Before you move on}]
What does \bn{বট} mean? (\textquotedblleft{}A banyan tree\textquotedblright{}.) Say it once more.
\end{tcolorbox}
\section[\texorpdfstring{\bn{ধান} (dhān)}{ধান (dhān)}]{\bn{ধান} (dhān) — paddy, rice plants}

\label{lesson:BN-C171-dhan}



\begin{quote}
Before the new one: say the Bengali for a rose, then the Bengali for a banyan tree.
\end{quote}

\subsection*{You'll want to know: \bn{ধান}}
\textbf{\bn{ধান}} — \emph{dhān} — \textquotedblleft{}paddy, rice plants\textquotedblright{}.

\begin{cousinweb}
\textbf{\bn{ধান}} (\emph{dhān}) is a tadbhava, a word inherited from Sanskrit and reshaped by sound change. It comes from \emph{dhānya}, \textquotedblleft{}grain; rice\textquotedblright{} (CDIAL 6778). CDIAL lists Bengali \emph{dhān}, \textquotedblleft{}growing or unhusked rice\textquotedblright{}.
\end{cousinweb}

\begin{grammarlens}[title={What you've built}]
One more word for this chapter.
\end{grammarlens}

\subsection*{Your turn}
\begin{itemize}
\raggedright
  \item \emph{Say it:} \emph{dhān}
  \item \emph{Say it:} \emph{dhān}, once more
  \item \emph{Say it:} say \emph{dhān}
  \item \emph{Recall:} say the Bengali for a rose, then the Bengali for a banyan tree, then say \emph{dhān} again
\end{itemize}

\begin{tcolorbox}[breakable,colback=teal!4,colframe=teal!35!black,title={Before you move on}]
What does \bn{ধান} mean? (\textquotedblleft{}Paddy, rice plants\textquotedblright{}.) Say it once more.
\end{tcolorbox}
\section[\texorpdfstring{\bn{জোয়ার} (joyār)}{জোয়ার (joyār)}]{\bn{জোয়ার} (joyār) — high tide}

\label{lesson:BN-C171-joyar}



\begin{quote}
Before the new one: say the Bengali for a banyan tree, then the Bengali for paddy.
\end{quote}

\subsection*{You'll want to know: \bn{জোয়ার}}
\textbf{\bn{জোয়ার}} — \emph{joyār} — \textquotedblleft{}high tide\textquotedblright{}.

\begin{cousinweb}
\textbf{\bn{জোয়ার}} (\emph{joyār}) is a tadbhava, a word inherited and reshaped by sound change. It comes from \emph{*javacāra}, \textquotedblleft{}flood-tide\textquotedblright{}, a form CDIAL reconstructs from its descendants (CDIAL 5167).
\end{cousinweb}

\begin{grammarlens}[title={What you've built}]
One more word for this chapter.
\end{grammarlens}

\subsection*{Your turn}
\begin{itemize}
\raggedright
  \item \emph{Say it:} \emph{joyār}
  \item \emph{Say it:} \emph{joyār}, once more
  \item \emph{Say it:} say \emph{joyār}
  \item \emph{Recall:} say the Bengali for a banyan tree, then the Bengali for paddy, then say \emph{joyār} again
\end{itemize}

\begin{tcolorbox}[breakable,colback=teal!4,colframe=teal!35!black,title={Before you move on}]
What does \bn{জোয়ার} mean? (\textquotedblleft{}High tide\textquotedblright{}.) Say it once more.
\end{tcolorbox}
